% ======================================================================
%  preamble.tex  --  shared preamble for the internship report
%  Separation of Concerns in Gaussian-Process Bayesian Optimization
% ======================================================================

% ---- Encoding, fonts, language -----------------------------------------
\usepackage[utf8]{inputenc}
\usepackage[T1]{fontenc}
\usepackage[english]{babel}
\usepackage[expansion=false]{microtype}

% ---- Math ---------------------------------------------------------------
\usepackage{amsmath,amssymb,amsthm,mathtools}
\usepackage{bm}

% ---- Layout -------------------------------------------------------------
\usepackage[a4paper,margin=2.6cm]{geometry}
\usepackage{setspace}
\onehalfspacing
\usepackage{parskip}
\usepackage{enumitem}
\setlist{itemsep=2pt,topsep=3pt,parsep=0pt}

% ---- Graphics, colour, tables ------------------------------------------
\usepackage{placeins}
\usepackage{graphicx}
\usepackage{xcolor}
\usepackage{booktabs}
\usepackage{array}
\usepackage{tabularx}
\usepackage{caption}
\usepackage{subcaption}
\captionsetup{font=small,labelfont=bf}

% ---- Headers ------------------------------------------------------------
\usepackage{fancyhdr}
\pagestyle{fancy}
\fancyhf{}
\fancyhead[L]{\small\itshape Separation of Concerns in GP Bayesian Optimization}
\fancyhead[R]{\small\thepage}
\renewcommand{\headrulewidth}{0.4pt}

% ---- Bibliography (natbib + bibtex) ------------------------------------
\usepackage[numbers,square,sort&compress]{natbib}

% ---- Hyperlinks (load late) --------------------------------------------
\usepackage[hidelinks]{hyperref}
\usepackage{cleveref}

% ======================================================================
%  Colour palette for provenance tags
% ======================================================================
\definecolor{colLit}{HTML}{4C6B8A}      % established literature
\definecolor{colSpec}{HTML}{2E8B8B}     % literature specialized to our setting
\definecolor{colContrib}{HTML}{B5651D}  % personal theoretical contribution
\definecolor{colExp}{HTML}{7A5CA5}      % personal experimental contribution
\definecolor{colSynth}{HTML}{3F7A3F}    % interpretation / synthesis
\definecolor{colOpen}{HTML}{9A2B2B}     % open question / conjecture

% Inline provenance badge --------------------------------------------------
\newcommand{\provbadge}[2]{%
  \hfill\raisebox{0.0ex}{\colorbox{#1!16}{\footnotesize\textcolor{#1!75!black}{\textsf{#2}}}}%
}
\newcommand{\Plit}{\provbadge{colLit}{established literature}}
\newcommand{\Pspec}{\provbadge{colSpec}{literature, specialized here}}
\newcommand{\Pcontrib}{\provbadge{colContrib}{personal contribution}}
\newcommand{\Pexp}{\provbadge{colExp}{personal experiment}}
\newcommand{\Psynth}{\provbadge{colSynth}{interpretation / synthesis}}
\newcommand{\Popen}{\provbadge{colOpen}{open question}}

% Standalone legend chip (for the introduction legend) --------------------
\newcommand{\legendchip}[2]{\colorbox{#1!16}{\footnotesize\textcolor{#1!75!black}{\textsf{#2}}}}

% ======================================================================
%  Theorem environments (numbered within section)
% ======================================================================
\theoremstyle{plain}
\newtheorem{theorem}{Theorem}[section]
\newtheorem{proposition}[theorem]{Proposition}
\newtheorem{lemma}[theorem]{Lemma}
\newtheorem{corollary}[theorem]{Corollary}
\theoremstyle{definition}
\newtheorem{definition}[theorem]{Definition}
\newtheorem{assumption}{Assumption}
\theoremstyle{remark}
\newtheorem{remark}[theorem]{Remark}
\newtheorem*{takeaway}{Takeaway}

% ======================================================================
%  Figure inclusion with automatic placeholder
%  \figauto[width]{path}{caption}{label}
%  If the PNG exists it is included; otherwise a labelled grey box is drawn
%  so the document always compiles.  Drop the real files into figures/ .
% ======================================================================
\newcommand{\figph}[2][\linewidth]{%
  \IfFileExists{#2}%
    {\includegraphics[width=#1,height=0.30\textheight,keepaspectratio]{#2}}%
    {\setlength{\fboxsep}{4pt}%
     \fbox{\parbox[c][0.24\textheight][c]{0.90\linewidth}{%
       \centering\ttfamily\footnotesize
       [\,figure placeholder\,]\\[4pt]\detokenize{#2}}}}%
}

\newcommand{\figauto}[4][\linewidth]{%
  \begin{figure}[htbp]
    \centering
    \figph[#1]{#2}
    \caption{#3}
    \label{#4}
  \end{figure}}

% ======================================================================
%  TODO / verification markers
% ======================================================================
\usepackage{todonotes}
\setuptodonotes{inline,size=\footnotesize}
\newcommand{\TODO}[1]{\todo[color=yellow!45,inline]{\textbf{TODO:} #1}}
% lightweight inline placeholder (safe inside titlepage / tabular / captions)
\newcommand{\TODOtext}[1]{{\colorbox{yellow!45}{\footnotesize\textbf{[TODO:} #1\textbf{]}}}}
\newcommand{\CITEQ}[1]{\textcolor{colOpen}{\textbf{[cite?]}}\marginpar{\footnotesize\textcolor{colOpen}{#1}}}

% ======================================================================
%  Notation macros
% ======================================================================
\newcommand{\R}{\mathbb{R}}
\newcommand{\N}{\mathbb{N}}
\newcommand{\E}{\mathbb{E}}
\newcommand{\Prob}{\mathbb{P}}
\newcommand{\X}{\mathcal{X}}
\newcommand{\Hk}{\mathcal{H}_k}
\newcommand{\GP}{\mathcal{GP}}
\newcommand{\inner}[2]{\langle #1, #2\rangle}
\newcommand{\norm}[1]{\left\lVert #1\right\rVert}
\newcommand{\abs}[1]{\left\lvert #1\right\rvert}
\newcommand{\Hnorm}[1]{\norm{#1}_{\Hk}}
\newcommand{\dist}{\operatorname{dist}}
\newcommand{\argmax}{\operatorname*{arg\,max}}
\newcommand{\argmin}{\operatorname*{arg\,min}}
\DeclareMathOperator{\vol}{vol}
\DeclareMathOperator{\HV}{HV}
\newcommand{\Fvec}{\bm{F}}
\newcommand{\prm}{\rho}                 % task transformation
\newcommand{\Rregret}{R}
\newcommand{\gain}{\gamma}
\newcommand{\rkhs}{\mathcal{H}}
\newcommand{\eps}{\varepsilon}
\newcommand{\epsinf}{\eps_{\infty}}
\newcommand{\RIGHT}{\textsc{Right}}
\newcommand{\WRONG}{\textsc{Wrong}}
\newcommand{\SoC}{\textsc{SoC}}
